\section{A fixed six-clock inverse for convex twist suspensions}
\label{sec:v79-suspensions}

We next verify the generating-action hypotheses on a class in which the
unknown is a function of two variables.  In particular the argument is
not restricted to kinetic energy plus a potential.  The observed clock
will be the return time of an explicitly specified contact suspension.
This construction distinguishes a genuine clock process from a device
that is instructed to evaluate an unknown stationary action.

Fix positive constants $\lambda,\Lambda,b,A_0,A_1$.  Consider all
$L\in C^\infty(\R^2)$ satisfying
\begin{equation}\label{eq:v79-convex-class}
 \lambda I_2\preceq D^2L(x,y)\preceq\Lambda I_2,
 \qquad -L_{xy}(x,y)\ge b,
 \qquad |L(0,0)|\le A_0,\quad |DL(0,0)|\le A_1.
\end{equation}
The class is assumed nonempty; the inequalities themselves constrain
the possible constants.  Bounds for higher derivatives will be imposed
only when quantitative stability is asserted.  The scalar configuration
coordinate and its canonical covector are fixed coordinates of this
mechanical experiment, unlike the unknown Euclidean embeddings in the
billiard application.

\begin{lemma}[Intrinsic determinism, convex composition, and twist]
\label{lem:v79-convex-twist}
The equations
\[
             p=-L_x(x,y),\qquad q=L_y(x,y)
\]
define a global smooth exact symplectic diffeomorphism
$F_L:(x,p)\mapsto(y,q)$ of $\R^2$.  For $n\ge1$ the function
\begin{equation}\label{eq:v79-action-min}
 S_n(s,t)=\min_{x_1,\ldots,x_{n-1}}
             \sum_{i=0}^{n-1}L(x_i,x_{i+1}),
             \qquad x_0=s,\quad x_n=t,
\end{equation}
is a unique smooth stationary action and is the generating function
of $F_L^n$.  Its mixed derivative satisfies
\begin{equation}\label{eq:v79-twist-bound}
 |(S_n)_{st}|\ge u_n:=b\left(\frac{b}{2\Lambda}\right)^{n-1}.
\end{equation}
At every match between $S_n$ and $S_{n+1}$ the stationary Schur
complement is positive and at most $2\Lambda$.  Equality of initial
covectors identifies the entire physical prefix, and the matching
function has derivative at most $-u_n^2/(2\Lambda)$ at each zero.
\end{lemma}
\begin{proof}
For fixed $x$, the derivative of $y\mapsto-L_x(x,y)$ is at least
$b$.  This map is strictly increasing and onto $\R$, since its values
grow at least linearly at both ends.  It has a unique smooth inverse.
For fixed $y$, the map $x\mapsto L_y(x,y)$ is strictly decreasing,
with derivative at most $-b$, and is likewise onto.  These two inverses
construct $F_L$ and its smooth inverse.  With $\theta=p\,dx$ and
$\ell_L(x,p)=L(x,y(x,p))$, the first variation gives
\begin{equation}\label{eq:v79-exact-primitive}
                       F_L^*\theta-\theta=d\ell_L.
\end{equation}
Taking an exterior derivative proves symplecticity.

The second variation of the sum in~\eqref{eq:v79-action-min}, with
endpoints fixed, is bounded below by
\[
 \lambda\sum_{i=0}^{n-1}(h_i^2+h_{i+1}^2)
                    =2\lambda\sum_{i=1}^{n-1}h_i^2.
\]
Coercivity and strict convexity give a unique minimum.  Its Hessian is
tridiagonal, with diagonal entries at most $2\Lambda$ and neighboring
entries $L_{xy}\le-b$.  The implicit function theorem gives smooth
dependence on $(s,t)$.  The corner-cofactor formula, including the
empty determinant for $n=1$, gives
\[
 |(S_n)_{st}|=
     \frac{\prod_{i=0}^{n-1}|L_{xy}(x_i,x_{i+1})|}
                  {\det H_{n-1}}
       \ge \frac{b^n}{(2\Lambda)^{n-1}},
\]
where the determinant bound is Hadamard's inequality for the positive
Hessian.  The Euler--Lagrange equations identify consecutive covectors,
so this minimizer follows $F_L$.  Conversely an orbit with the given
endpoints is stationary and hence equals the unique minimizer.
This also proves the stated generating-function interpretation.

A prefix match is therefore a match of actual trajectories, not merely
of stationary-value continuations.  The last-site Schur complement of
the full positive Hessian is positive and no greater than that site's
diagonal entry $2\Lambda$.  Differentiating stationary composition gives
$H_s=-(S_n)_{st}^2/B$, exactly as in~\eqref{eq:v78-root-sign}.
The lower bound for the twist proves the final assertion.
\end{proof}

\subsection{An elapsed-time realization}

The elementary contact construction below is part of the classical
relation between exact symplectic action and return time; compare the
contact realization of action functions in~\cite{Hutchings2016}.
We include the construction to specify the observation process and its
clock, not to claim a new construction of contact suspensions.

\begin{proposition}[Positive action as a return-time roof]
\label{prop:v79-contact}
Let $F$ be a global exact symplectic diffeomorphism of $(\R^2,d\theta)$,
$\theta=p\,dx$, with $F^*\theta-\theta=d\ell$.  Suppose
$r=c+\ell\ge r_0>0$.  The quotient of $\R^2\times\R$ by
\begin{equation}\label{eq:v79-contact-deck}
       G(x,p,z)=(F(x,p),z-r(x,p))
\end{equation}
is a smooth contact manifold with contact form induced by
$\alpha=dz+\theta$.  Its complete Reeb flow has a global section
$z=0$, first-return map $F$, and first-return time $r$.
For $F=F_L$ in Lemma~\ref{lem:v79-convex-twist}, the time of its
$n$th return from $(s,p)$ with terminal coordinate $t$ is
\begin{equation}\label{eq:v79-return-time}
                            nc+S_n(s,t).
\end{equation}
\end{proposition}
\begin{proof}
The identity $G^*(dz+\theta)=dz-dr+F^*\theta=dz+\theta$ proves
invariance.  Also $\alpha\wedge d\alpha=dz\wedge dp\wedge dx$
never vanishes.  The integer action is free and properly discontinuous:
under $G^k$, $k>0$, the $z$ coordinate decreases by at least $kr_0$;
under negative iterates it increases by at least $|k|r_0$.  Thus only
finitely many iterates can meet two fixed compact subsets.  The quotient
is a smooth manifold and the invariant form descends.

On the covering space the Reeb field is $\partial_z$, since
$\alpha(\partial_z)=1$ and $d\alpha(\partial_z,\cdot)=0$.
It is invariant under $G$ and its translation flow is complete.
Starting at $(x,p,0)$, the first time its class returns to the section
is $r(x,p)$, at which
$(x,p,r(x,p))$ is identified with $(F(x,p),0)$.  Positivity of the
roof excludes any earlier section crossing; negative powers correspond
to past crossings.  The section is an embedded copy of $\R^2$, since
two of its points cannot be related by a nonzero iterate.
Summing successive roofs gives
$nc+\sum_{i=0}^{n-1}L(x_i,x_{i+1})$ in the convex twist case.
The orbit is the unique minimizer for its endpoints by the preceding
lemma, proving~\eqref{eq:v79-return-time}.
\end{proof}

The clock in this proposition is elapsed Reeb time on the stated
contact manifold.  It is not asserted to be ordinary elapsed time of
an unsuspended discrete mechanical map.  A choice of the primitive's
additive constant changes the roof, and is therefore physically visible
to this observation.  The known deadlines fix that constant.

\begin{theorem}[Class-wide six-deadline reconstruction]
\label{thm:v79-six-clock}
Fix a target radius $R>0$, a finite integer $n\ge1$, and the constants
in~\eqref{eq:v79-convex-class}.  There are a fixed initial phase box,
fixed endpoint gates, a fixed release-delay interval, and three fixed
deadlines for each of the return counts $n$ and $n+1$, all independent
of $L$, with the following property.

In the contact suspension of $F_L$, launch a state on the section after
an independent uniform external release delay.  At the specified return
count retain the record only if it is in the endpoint gate and has
arrived by the specified deadline, with a positive endpoint-dependent
efficiency.  The source density and efficiency may be unknown and may
differ between the two return counts, but are shared by the three
deadlines of each count.  Record only the initial and terminal scalar
positions upon retention, and a failure symbol otherwise.  The six
conditional endpoint laws determine the full function
$L|_{[-R,R]^2}$, including its additive constant.

No root-selection mark, observed covector, individual return time,
or action evaluation is part of the retained data.  On subclasses with
bounded finite-order derivatives on the compact phase and configuration
domains constructed below, and positive source and efficiency floors,
reconstruction is uniformly locally Lipschitz from $C^{k+2}$ laws to
$C^k([-R,R]^2)$ for every finite $k$.  For clock-dependent log-weight
errors of $C^{k+2}$ size $\epsilon$, the same reconstruction has
additional error at most $C_k\epsilon$.
\end{theorem}
\begin{proof}
We construct all settings from the class constants.  The stationarity
equation at a configuration value $x$ is
\[
                    L_y(x_-,x)+L_x(x,x_+)=0.
\]
For fixed $x,x_+$ its left side is strictly decreasing in $x_-$ with
derivative at most $-b$ and has one zero.  At $x_-=0$ its absolute
value is at most $2A_1+2\Lambda|x|+\Lambda|x_+|$, using the
Hessian bound and the bound on $DL(0,0)$.  Consequently
\begin{equation}\label{eq:v79-backward-bound}
 |x_-|\le\frac{2A_1+2\Lambda|x|+\Lambda|x_+|}{b}.
\end{equation}
Set $r_0=r_1=R+1$ and
\[
 r_{j+1}=(2\Lambda/b)r_j+(\Lambda/b)r_{j-1}+2A_1/b,
 \qquad S=r_{n+1}+1,
 \qquad I=[-S,S],\quad J_0=[-R-1,R+1].
\]
The ratios in this recurrence are nonnegative and $\Lambda\ge b$,
so its bounds increase.  Starting from any final pair
$(t,v)\in J_0^2$, backward solution through $n$ sites gives an
initial coordinate in the interior of $I$, with a unit buffer.
Its whole path is stationary; strict convexity identifies it with
both actions $S_n(s,t)$ and $S_{n+1}(s,v)$ on the fixed gates
$I\times J_0$.  It therefore supplies a matching root for every
$(t,v)\in J_0^2$.  Lemma~\ref{lem:v79-convex-twist} and the downward-zero
lemma show that this root is unique, without specifying it to the observer.

Put
\[
 C_0=A_0+\frac{A_1^2}{2\lambda},\qquad
 M=A_0+\sqrt2 A_1S+\Lambda S^2.
\]
The quadratic lower bound for $L$ gives $L\ge-C_0$ globally.
A straight interpolating configuration path for either endpoint gate
stays in $[-S,S]$ and has each term at most $M$.  Thus for
$q\in\{n,n+1\}$ both actions satisfy the fixed bounds
\[
       w_-=-(n+1)C_0\le S_q\le w_+=(n+1)M.
\]
Take $c=C_0+1$ in Proposition~\ref{prop:v79-contact}, so every roof
is at least one.  Set
\begin{equation}\label{eq:v79-deadlines}
 T_j=w_++j,\quad A=w_+-w_-+4,\quad
 \mathsf D_{q,j}=qc+T_j,\qquad j=1,2,3.
\end{equation}
Then $1\le T_j-S_q\le A-1$ on the fixed gates.

The source phase box can be fixed as well.  The inequality
\[
 L(x,y)\ge\frac{\lambda}{4}(x^2+y^2)-C_1,
       \qquad C_1=A_0+A_1^2/\lambda,
\]
follows by absorbing the linear term in the quadratic lower bound.
At a minimizing path of length $q\le n+1$ it implies
\[
 \sum_{i=1}^{q-1}x_i^2
       \le\frac{2q}{\lambda}(M+C_1).
\]
Hence all its coordinates lie in $[-K,K]$ for
\[
 K=\max\left\{S,
          \sqrt{2(n+1)(M+C_1)/\lambda}\right\}.
\]
The initial covector has magnitude at most $A_1+\sqrt2\Lambda K$.
Choose a strictly larger fixed phase box and a normalized smooth source
positive on the required compact subset.  Such a source can even be
fixed across the entire class; the theorem allows any unknown source
with the stated positivity and sharing properties.  Different full
source laws are not confused with different acquisition rules.

Let $\alpha_0$ be the external release delay, uniform on $(0,A)$.
Stopping at return $q$ and comparing the clock with $\mathsf D_{q,j}$
requires only counting section crossings and reading a clock.
By~\eqref{eq:v79-return-time}, the event is exactly
\[
           \alpha_0+qc+S_q(s,t)<\mathsf D_{q,j}
               \quad\Longleftrightarrow\quad
                    \alpha_0<T_j-S_q(s,t).
\]
The delay is external: it is not an assertion that a long interval
before a section crossing contains no earlier crossings.  Changing
from the initial $p$ to the terminal $t$ has the nonzero Jacobian
$|(S_q)_{st}|$.  Integrating the independent delay therefore yields
unnormalized density
\[
 \frac1A e_q(s,t)\rho_q(s,-(S_q)_s(s,t))
             |(S_q)_{st}(s,t)|\{T_j-S_q(s,t)\}.
\]
Normalization gives precisely the shared-weight model, separately
for $q=n$ and $q=n+1$.  The observer is not instructed to evaluate
$S_q$; this expression is the mathematical derivation of a clock event
in the specified flow.

Recover $S_n,S_{n+1}$ and their absolute constants from their triples.
At the unique match subtract the prefix to obtain $L(t,v)$ on all of
$J_0^2$, by Theorem~\ref{thm:v78-generating}.  This recovers the
required compact square.  For quantitative stability, the positive
interior Hessian has inverse norm at most $(2\lambda)^{-1}$;
implicit differentiation bounds all required action derivatives on
the fixed compact region.  The twist estimate and the Schur bound
control the matching derivative.  The unit initial buffer, restriction
to the smaller target square, and these derivative bounds give uniform
root collars and sign margins.  For example, on a collar where
$|H_s|\ge u_n^2/(4\Lambda)$, the value of $H$ at its two endpoints
has the corresponding fixed positive magnitude.  Bounded next
derivatives make a common collar possible.  The action-anchor bound
in Proposition~\ref{prop:v78-anchor-bound} supplies a usable anchor
among finitely many corners.  Composing the clock inverse with the
root and subtraction estimates proves the $C^{k+2}$ to $C^k$ bound.
Finally Theorem~\ref{thm:v79-nuisance} gives a $C^{k+2}$ action bias
at most $C\epsilon$; the same root estimate propagates it to $L$.
\end{proof}

\begin{remark}[Breadth and the retained mechanical subfamily]
\label{rem:v79-breadth}
For example
\[
 L_*(x,y)=\frac\mu2(y-x)^2+\frac\kappa2(x^2+y^2)
\]
has Hessian eigenvalues $\kappa$ and $2\mu+\kappa$ and mixed
entry $-\mu$.  Adding a sufficiently small, nonseparable smooth bump
supported in a compact subset of $\R^2$ preserves strict versions of
\eqref{eq:v79-convex-class}.  The bump may be flat at its support
boundary and hence nonanalytic.  The inverse recovers this entire
unknown two-variable perturbation on a prescribed square, not merely
a finite jet or a parameter vector.

The original family $L_{\mu,P}(x,y)=\mu(y-x)^2/2+P(x)$ with
\eqref{eq:v78-mechanical-class} is also contained in a class of the
form~\eqref{eq:v79-convex-class}.  Indeed its determinant is
$\mu P''$, its trace is $2\mu+P''$, and its mixed entry is $-\mu$;
one may take
\[
 \lambda=\frac{\mu_-\kappa}{2\mu_++\Lambda},\qquad
 \Lambda_{\rm joint}=2\mu_++\Lambda,\qquad b=\mu_-.
\]
Here the two occurrences of $\Lambda$ in the fraction refer to the
potential-curvature bound in~\eqref{eq:v78-mechanical-class}.
The mass and potential formulas, their sharper elementary bounds,
and their full proofs are retained in the next section.  For that
family the suspension construction supplies a return-clock realization
of its action-threshold experiment.  It does not change the time
interpretation of the unsuspended mechanical map.
\end{remark}
